%% kapa: exam=1 round=19 year=2008 period=1 form=A subject=accounting subject_ko=회계학 seq=40 number=120
%% source: Q-Net 감정평가사 기출문제 (cst00309) artlSeq=1001962 file=extracted/1008/1001962/2008년 제19회 감정평가사 1차시험 기출문제/1교시.hwp
%% conversion: kordoc 4.24.51 -> tools/gen1.py (auto); review=proofread-sonnet; flags=none
\begin{question}{120}
다음 항목들은 재무제표 확정 전 대차대조표일 후 발생한 사건의 사례이다. 이 중 재무제표를 수정할 필요가 없는 것은?
\begin{choices}
  \item 대차대조표일 현재 이미 자산의 가치가 하락되었음을 나타내는 정보를 대차대조표일 후에 입수하는 경우
  \item 건설 중인 교량이 천재지변으로 인하여 대차대조표일과 재무제표가 사실상 확정된 날 사이에 붕괴된 경우
  \item 대차대조표일 이전에 구입한 자산의 취득원가 또는 매각한 자산의 금액을 대차대조표일 후에 결정하는 경우
  \item 전기 또는 그 이전기간에 발생한 회계적 오류를 대차대조표일 후에 발견하는 경우
  \item 대차대조표일 이전에 존재하였던 소송사건의 결과가 대차대조표일 후에 확정되어 이미 인식한 손실금액을 수정해야 하는 경우
\end{choices}
\end{question}
\answer{2}
